\section{Estimation, testing and model checking}\label{sec:estimation}

This section turns the model of \cref{sec:framework} into inference. Everything rests on
the log-likelihood \eqref{eq:th-loglik}. \Cref{sec:est-mle} evaluates it exactly for
exponential kernels and discusses the background. \Cref{sec:est-em} maximises it by EM
over the latent branching structure, \cref{sec:est-bayes} gives a Bayesian version with
shrinkage on $\Amat$, and \cref{sec:est-inar} treats binned counts. We then test whether
one stream excites another (\cref{sec:est-lrt}), check goodness of fit
(\cref{sec:est-gof}), and calibrate shock effects against placebo calendars
(\cref{sec:est-placebo}). Most ingredients are known, and we cite them as we go. To our
knowledge, two results are new in this form:
\begin{itemize}
  \item with the observation window kept exactly in the complete-data likelihood, the
    EM update of an exponential decay rate is a strictly concave one-dimensional problem
    with a moment-matching reading (\cref{prop:est-mstep});
  \item an exact formula for the attenuation of contagion estimates computed from
    binned counts (\cref{prop:est-binning}).
\end{itemize}
We also make explicit an elementary fact about M-matrices, which turns the stationarity
constraint into a truncated-gamma Gibbs update (\cref{lem:est-trunc}). Proofs, the Gibbs
sampler, nonparametric kernels, the conditions and proofs of the classical boundary and
time-rescaling results, and the Davies problem are in \cref{app:est-supp}; computational
details are in \cref{app:computation}.

Throughout, $\vartheta$ denotes the full parameter vector: the background coefficients
$\theta_k$, the shock loadings $\upsilon_k$, the response parameters, $\Amat$ and the kernel
parameters. We reserve $\theta_k$ for the background, as in \cref{ass:th-model}. We write
$n_k=N_k(T)$ and $n_{\mathrm{ev}}=\sum_kn_k$ for the total number of incidents; $n$
without a subscript is the branching ratio, as in \cref{sec:framework}.

\subsection{Maximum likelihood}\label{sec:est-mle}

\paragraph{Separation by target stream.} The log-likelihood \eqref{eq:th-loglik} is a sum
$\ell_T=\sum_k\ell_{T,k}$ of stream terms. The term $\ell_{T,k}$ depends only on the
parameters of $\lambda_k$: on $\theta_k$, on $\upsilon_k$ and $\psi_k$, and on row $k$ of
$\Amat$ with its kernels. Unless a parameter is shared across rows, such as a common decay
rate, the likelihood therefore separates. Each row of the model is fitted to its own
stream, and the other streams enter only as history. We use this separation for testing
(\cref{sec:est-lrt}) and for parallel computation (\cref{app:computation}).

\paragraph{Exponential kernels.} Let $g_{kj}(u)=\beta_{kj}e^{-\beta_{kj}u}$ and
$\psi_k(u)=\eta_ke^{-\eta_ku}$ for $u>0$, and let $\Psi_k(u)=\int_0^u\psi_k$, with
$\Psi_k(u)=0$ for $u\le0$. With the decay sums
$R_{kj}(t)=\sum_{l:\,t^j_l<t}e^{-\beta_{kj}(t-t^j_l)}$, the intensity and the compensator
on $[0,T]$ are
\begin{align}
  \lambda_k(t)&=\mu_k(t)+\upsilon_k\sum_{e:\,\tau_e<t}S_e\,\eta_ke^{-\eta_k(t-\tau_e)}
  +\sum_{j=1}^M\alpha_{kj}\beta_{kj}R_{kj}(t),\label{eq:est-lambda}\\
  \Lam_k(T)&=\int_0^T\mu_k(t)\,\dd t+\upsilon_k\sum_eS_e\Psi_k^T(e)
  +\sum_{j=1}^M\alpha_{kj}G_{kj},\label{eq:est-comp}
\end{align}
where $\Psi^T_k(e)=\Psi_k(T-\tau_e)-\Psi_k(-\tau_e)$ is the share of the response to
event $e$ that falls inside the window, and
$G_{kj}=\sum_{l=1}^{n_j}\{1-e^{-\beta_{kj}(T-t^j_l)}\}$ is the expected number of in-window
offspring per unit of $\alpha_{kj}$. For a pair $(k,j)$, let $s_1<s_2<\cdots$ be the
merged event times of streams $j$ and $k$, with $d_m$ the stream of $s_m$. The exponential
kernel is memoryless, so
\begin{equation}\label{eq:est-recursion}
  R_{kj}(s_m)=e^{-\beta_{kj}(s_m-s_{m-1})}\big\{R_{kj}(s_{m-1})+\ind{d_{m-1}=j}\big\},
  \qquad R_{kj}(s_1)=0 .
\end{equation}
For each pair, \eqref{eq:est-recursion} costs $O(n_j+n_k)$, so all intensities at the
events cost $O(Mn_{\mathrm{ev}})$ operations, against $O(n_{\mathrm{ev}}^2)$ for direct
summation. The recursion is due to \citet{ozaki1979}; \cref{app:computation} gives the
analogous recursions for the gradient and the EM statistics.

\paragraph{Asymptotic theory.} \Citet{ogata1978} proved consistency, asymptotic normality
and efficiency of the maximum likelihood estimator for stationary point processes. His
conditions are stationarity and ergodicity, identifiability, a compact parameter space
with $\vartheta_0$ in its interior, and smoothness and moment conditions on the intensity
and its first three derivatives. Under them,
$\sqrt T(\hat\vartheta-\vartheta_0)\Rightarrow N(0,\bm I^{-1})$, where
$\bm I=\sum_k\E\{\lambda_k(0)^{-1}\nabla\lambda_k(0)\nabla\lambda_k(0)^{\top}\}$ is the Fisher
information per unit time under the stationary law. The shock-augmented model with a
spline background is not stationary. A fixed-dimensional theory applies if the number
$K_b$ of basis functions is held fixed as $T$ grows. Our pre-registered design instead
fixes the knot spacing $h_b$ (\cref{sec:plan-models}), so $K_b$ grows in proportion to $T$
and the background is a semiparametric nuisance, whose theory we do not develop. We use the
normal approximation with the observed information (\cref{app:computation}) away from the
boundary, and parametric bootstrap or posterior draws where it is in doubt: near
$\alpha_{kj}=0$ (\cref{sec:est-lrt}), near criticality, and for nonlinear functionals such
as the surge decomposition (\cref{sec:at-uncertainty}).

\paragraph{Constraints.} The parameters satisfy $\alpha_{kj}\ge0$, $\upsilon_k\ge0$ and
$\beta_{kj},\eta_k>0$, so $\lambda_k\ge\mu_k>0$ and the likelihood is defined for every
$\Amat\ge\zero$. The condition $\spr(\Amat)<1$ is needed only for the stationary
interpretation, for $\Mmat$ and for the surge decomposition. We fit without it and report
$\spr(\hat\Amat)$ with its uncertainty, reading an estimate near or above one as a symptom
of misspecification (\cref{thm:spurious}) rather than as evidence of an explosive process.
If the constraint is imposed, the M-step for $\Amat$ with the decays fixed stays convex,
because $\log\spr(\Amat)$ is convex in the log-entries \citep{kingman1961}; in the Bayesian
fit we impose it by truncation (\cref{lem:est-trunc}).

\paragraph{The background, and why its flexibility matters.} The background is
$\mu_k(t)=\exp\{B(t)^{\top}\theta_k\}$. Here $B$ consists of cubic B-splines in calendar
time with $K_b$ basis functions, plus day-of-week and holiday indicators
(\cref{sec:plan-models}). We maximise the penalised log-likelihood
\begin{equation}\label{eq:est-pen}
  \ell^{\mathrm{pen}}_T(\vartheta)=\ell_T(\vartheta)-\tfrac12\sum_{k}\lambda^{\mathrm{pen}}_k\,
  \theta_k^{\top}P\,\theta_k .
\end{equation}
Two choices of $P$ are natural: the second-difference penalty on adjacent spline
coefficients \citep{eilers1996}, which leaves linear trends unpenalised, and
$P=\Imat-\mathbf 1\mathbf 1^{\top}/K_b$, which shrinks towards a constant background. Our
implementation offers the second, and by default uses no penalty, with the flexibility set
by the knot spacing alone. The background is not a nuisance to be fitted as flexibly as
possible. It is part of the identification strategy, and it can fail in two directions.
\begin{itemize}
  \item \emph{Too rigid.} Rate variation that the background cannot represent is absorbed
    by the excitation terms. \Cref{thm:spurious} makes this exact for a constant
    background. A process with no contagion but a fluctuating rate, such as one driven by
    unrecorded shocks or a regime-switching level, has second-order pseudo-true branching
    ratio $n^*=1-(1+\kappa)^{-1/2}>0$, and its likelihood pseudo-true value also has $n>0$.
  \item \emph{Too flexible.} A background that can vary on the time scale of a cascade can
    fit the realised clusters directly. In the limit the problem is degenerate. If $\mu_k$
    ranges over all nonnegative functions, then
    $\sup_{\mu}\ell_T(\mu,\Amat=\zero)=+\infty$: take $\mu$ equal to $(2\epsilon)^{-1}$ on
    $\epsilon$-neighbourhoods of the events and let $\epsilon\downarrow0$. The likelihood alone
    therefore cannot separate contagion from background. The separation comes only from
    restricting the roughness of $\mu_k$.
\end{itemize}
Identification thus rests on a separation of time scales: the background should be smooth
on the cascade time scale $1/\{\beta(1-n)\}$ of \cref{prop:th-spectrum}, while still
following seasonal and secular movement. For the primary analysis we fix the knot spacing
$h_b$ in advance at a multiple of the longest plausible cascade time scale, and report the
estimates for a pre-registered ladder of spacings (\cref{sec:plan-models,sec:plan-sens}). An
excitation estimate that vanishes as the background becomes more flexible, while the fit
stays adequate, is not robustly distinguishable from background variation. As a check, we
compare the pre-registered spacing with the one that maximises a blocked cross-validated log
score,
\begin{equation}\label{eq:est-cv}
  \mathrm{CV}=\sum_{I\in\mathcal I}\sum_{k=1}^{M}\Big[\sum_{t^k_i\in I}
  \log\hat\lambda^{(-I)}_k(t^k_i)-\int_{I}\hat\lambda^{(-I)}_k(t)\,\dd t\Big],
\end{equation}
over contiguous blocks $I$ much longer than the cascade time scale, where
$\hat\lambda^{(-I)}$ is fitted without block $I$ but conditions on the full history. In a
held-out block, background bumps fitted to clusters elsewhere do not help, whereas genuine
excitation still predicts clustering.

\subsection{EM via the latent branching structure}\label{sec:est-em}

By the cluster representation (\cref{sec:th-cluster}), every incident is a background
immigrant, a shock immigrant, or an offspring of an earlier incident. Let $Z^k_i$ be the
parent label of $t^k_i$. It takes the value $0$ (background), $e$ (shock), or $(j,l)$ with
$t^j_l<t^k_i$ (offspring of $t^j_l$). Given the labels, the immigrants and the offspring
of each incident form independent Poisson processes on $[0,T]$. The complete-data
log-likelihood is therefore
\begin{equation}\label{eq:est-complete}
\begin{aligned}
  \ell^{c}_T(\vartheta)={}&\sum_k\Big[\sum_{i:\,Z^k_i=0}\log\mu_k(t^k_i)-\int_0^T\mu_k\Big]
  +\sum_{k,e}\Big[\sum_{i:\,Z^k_i=e}\log\{\upsilon_kS_e\psi_k(t^k_i-\tau_e)\}-\upsilon_kS_e\Psi^T_k(e)\Big]\\
  &+\sum_{k,j}\Big[\sum_{(i,l):\,Z^k_i=(j,l)}\log\{\alpha_{kj}g_{kj}(t^k_i-t^j_l)\}
  -\alpha_{kj}\,\mathcal G_{kj}\Big],
  \qquad \mathcal G_{kj}=\sum_{l=1}^{n_j}\int_0^{T-t^j_l}g_{kj}(u)\,\dd u .
\end{aligned}
\end{equation}
The window terms $\Psi^T_k(e)$ and $\mathcal G_{kj}$ are the parts of the response and
offspring kernels that fall inside $[0,T]$. For exponential kernels, $\mathcal G_{kj}=G_{kj}$
of \eqref{eq:est-comp}. Summing $\exp(\ell^c_T)$ over all labellings recovers the
likelihood, because $\lambda_k(t^k_i)$ is the sum of the parent contributions. This is the
EM of \citet{veen2008} for the ETAS model, the stochastic declustering of
\citet{marsan2008} and the crime model of \citet{mohler2011}, extended here to exogenous
shocks.

\paragraph{E-step.} The posterior label probabilities are
\[
  p^{k,0}_i=\frac{\mu_k(t^k_i)}{\lambda_k(t^k_i)},\qquad
  p^{k,e}_i=\frac{\upsilon_kS_e\psi_k(t^k_i-\tau_e)}{\lambda_k(t^k_i)},\qquad
  p^{k,jl}_i=\frac{\alpha_{kj}g_{kj}(t^k_i-t^j_l)}{\lambda_k(t^k_i)} .
\]
The M-step needs only the aggregates
\begin{equation}\label{eq:est-stats}
\begin{gathered}
  N^0_k=\sum_ip^{k,0}_i,\qquad N_{ke}=\sum_ip^{k,e}_i,\qquad W_{ke}=\sum_ip^{k,e}_i(t^k_i-\tau_e),\\
  N_{kj}=\sum_{i,l}p^{k,jl}_i,\qquad W_{kj}=\sum_{i,l}p^{k,jl}_i(t^k_i-t^j_l).
\end{gathered}
\end{equation}
For exponential kernels,
$\sum_lp^{k,jl}_i=\alpha_{kj}\beta_{kj}R_{kj}(t^k_i)/\lambda_k(t^k_i)$ and
$\sum_lp^{k,jl}_i(t^k_i-t^j_l)=\alpha_{kj}\beta_{kj}R^{(1)}_{kj}(t^k_i)/\lambda_k(t^k_i)$, where
$R^{(1)}_{kj}(t)=\sum_{t^j_l<t}(t-t^j_l)e^{-\beta_{kj}(t-t^j_l)}$ obeys a recursion like
\eqref{eq:est-recursion} (\cref{app:computation}). The E-step therefore costs $O(Mn_{\mathrm{ev}})$
operations, and the $n_{\mathrm{ev}}\times n_{\mathrm{ev}}$ matrix of parent probabilities is never formed.

\begin{proposition}[M-step]\label{prop:est-mstep}
Fix the aggregates \eqref{eq:est-stats}, with $N_{kj}>0$ and $\sum_eN_{ke}>0$. Let
$Q(\vartheta)=\E\{\ell^c_T(\vartheta)\mid\text{data},\vartheta^{(r)}\}$, which is the
expression obtained from \eqref{eq:est-complete} by replacing label indicators with the
probabilities $p$.
\begin{enumerate}[(a)]
  \item \emph{Background.} For constant $\mu_k$, $Q$ is maximised at $\mu_k=N^0_k/T$. For
    $\mu_k=\exp(B^{\top}\theta_k)$, $\theta_k$ maximises the strictly concave function
    $\sum_ip^{k,0}_iB(t^k_i)^{\top}\theta_k-\int_0^Te^{B(t)^{\top}\theta_k}\dd t
    -\tfrac12\lambda^{\mathrm{pen}}_k\theta_k^{\top}P\theta_k$, provided the basis
    functions are linearly independent on $[0,T]$.
  \item \emph{Branching ratios and loadings.} For fixed kernel shapes,
    \[
      \alpha_{kj}=\frac{N_{kj}}{\mathcal G_{kj}},\qquad
      \upsilon_k=\frac{\sum_eN_{ke}}{\sum_eS_e\Psi^T_k(e)} .
    \]
  \item \emph{Decay rates.} For exponential $g_{kj}$, write $r_l=T-t^j_l$ (for the window
    with a burn-in segment, see \cref{rem:est-edge}) and
    \[
      m_j(\beta)=\frac{\sum_l\int_0^{r_l}u\,e^{-\beta u}\,\dd u}{\sum_l\int_0^{r_l}e^{-\beta u}\,\dd u}
      =\frac{\sum_l\{1-(1+\beta r_l)e^{-\beta r_l}\}}{\beta\sum_l\{1-e^{-\beta r_l}\}} .
    \]
    With $\alpha_{kj}$ profiled out, $Q$ is strictly concave in $\beta_{kj}$. The function
    $m_j$ decreases strictly from $m_j(0+)=\sum_lr_l^2/(2\sum_lr_l)$ to $0$. If
    $W_{kj}/N_{kj}<m_j(0+)$, the maximiser is the unique root of
    \[
      m_j(\beta_{kj})=W_{kj}/N_{kj}.
    \]
    Otherwise $Q$ increases as $\beta_{kj}\downarrow0$, and the update stops at a
    pre-set lower bound. In words, the mean lag of the in-window kernel equals the
    responsibility-weighted mean lag. A decay shared by a set $\mathcal P$ of pairs solves
    $\sum_{(k,j)\in\mathcal P}N_{kj}m_j(\beta)=\sum_{(k,j)\in\mathcal P}W_{kj}$. The response
    rate $\eta_k$ of an exponential $\psi_k$ solves the same equation, with $N_{kj},W_{kj}$
    replaced by $\sum_eN_{ke},\sum_eW_{ke}$, and the parent windows replaced by the shock
    windows with weights $S_e$.
  \item \emph{Monotonicity.} If $\vartheta^{(r+1)}$ maximises $Q$, or merely satisfies
    $Q(\vartheta^{(r+1)})\ge Q(\vartheta^{(r)})$, then
    $\ell_T(\vartheta^{(r+1)})\ge\ell_T(\vartheta^{(r)})$. With the penalty
    \eqref{eq:est-pen}, the same holds for $\ell^{\mathrm{pen}}_T$.
\end{enumerate}
\end{proposition}

The proof is in \cref{app:est-proofs}. The window terms $\mathcal G_{kj}$ and
$\Psi^T_k(e)$ matter. Dropping them gives the familiar edge-free updates
$\alpha_{kj}=N_{kj}/n_j$ and $\beta_{kj}=N_{kj}/W_{kj}$, an EM algorithm for a
pseudo-likelihood rather than for $\ell_T$. It underestimates $\alpha_{kj}$ and the loadings,
by a relative amount of about $1/(\beta_{kj}T)$, and overestimates decay rates. This is
negligible for a short kernel in a long window, but not for a salient event shortly before
$T$, which is exactly the case of real-time attribution. \Cref{rem:est-edge} in
\cref{app:est-edge} gives the details and the window terms for a burn-in segment.

\Cref{alg:em} in \cref{app:est-supp} summarises the iteration. Three practical points about \cref{alg:em} matter. First, the likelihood can be
multimodal in the decay rates; when excitation is weak, a spurious mode with tiny $\alpha$
and very large $\beta$ can appear, so our implementation starts EM from each local maximum
of a grid profile in the decays. When excitation is absent, the likelihood can also have a
ridge on which $\beta\to0$ while $\alpha$ grows and the expected number of in-window
offspring per incident stays small; on that ridge $\hat\alpha$ alone is not a usable
estimate of the branching ratio (\cref{sec:sim-spurious}). Second, EM converges linearly
and is slow when some $\alpha_{kj}$ is near $0$; SQUAREM extrapolation \citep{varadhan2008},
with steps kept only when they raise the likelihood, accelerates it, and null fits set
$\alpha_{kj}=0$ exactly. Third, the probabilities $p^{k,\cdot}_i$ give a stochastic
declustering of the realised incidents \citep{marsan2008}, used in \cref{rem:at-expost}.
The partial M-steps of \citet{veen2008} and the nonparametric kernels of \citet{lewis2011}
(\cref{sec:est-nonpar}) fit the same template.

\subsection{Bayesian estimation with shrinkage on \texorpdfstring{$\Amat$}{A}}\label{sec:est-bayes}

With many streams, $\Amat$ has $M^2$ entries, most of which we expect to be small. We use
priors that shrink small entries strongly and leave large ones nearly untouched: a
local--global gamma prior, whose marginal is a horseshoe-type prior placed directly on the
nonnegative $\alpha_{kj}$ \citep{carvalho2010,armagan2011}, or a spike-and-slab prior with
inclusion indicators $\iota_{kj}$ \citep{linderman2014}; \cref{app:est-gibbs} gives both.
Stationarity is imposed by truncating the prior to $\{\spr(\Amat)<1\}$. The next lemma, a
standard rank-one fact about nonsingular M-matrices, makes the truncation explicit one
entry at a time, in the form the sampler needs.

\begin{lemma}[The stationarity region, entry by entry]\label{lem:est-trunc}
Let $\Amat\ge\zero$, fix $(k,j)$, and let $\Amat^{(-kj)}$ be $\Amat$ with $\alpha_{kj}$ set
to $0$. Assume $\spr(\Amat^{(-kj)})<1$; this holds whenever $\spr(\Amat)<1$, by
Perron--Frobenius monotonicity. Let $\Mmat^{(-kj)}=(\Imat-\Amat^{(-kj)})^{-1}$. For
$x\ge0$,
\[
  \spr\big(\Amat^{(-kj)}+x\,\bm e_k\bm e_j^{\top}\big)<1
  \quad\Longleftrightarrow\quad
  x\,\Mmat^{(-kj)}_{jk}<1 .
\]
If $\Mmat^{(-kj)}_{jk}=0$, every $x\ge0$ is admissible.
\end{lemma}

The proof is in \cref{app:est-proofs}.

The threshold has a direct reading: adding the edge $j\to k$ with weight $x$ closes a
feedback loop with gain $x\,\Mmat^{(-kj)}_{jk}$, and the cascade stays subcritical exactly
when the gain is below one. An edge that closes no loop, such as an online-to-offline entry
in a model without offline-to-online feedback, is unconstrained.

\paragraph{Gibbs sampling by data augmentation.} Given the parent labels $Z$, the
complete-data likelihood \eqref{eq:est-complete} is a product of Poisson-process
likelihoods, so each parameter has a simple full conditional. In particular the window
term enters the branching ratios exactly, and
$\alpha_{kj}\mid\cdot\sim\operatorname{Gamma}(a_1+N_{kj},\,\omega_{kj}+G_{kj})$ truncated to
$[0,1/\Mmat^{(-kj)}_{jk})$, which is sampled by inversion (\cref{lem:est-trunc}).
\Cref{app:est-gibbs} gives the full sweep, including a block update of the spike-and-slab
indicators and a slice-sampling update of the decay rates. Each posterior draw of
$\vartheta$ gives draws of $\spr(\Amat)$, $\Mmat$ and the surge decomposition of
\cref{thm:attribution}, so parameter uncertainty propagates to every reported quantity
(\cref{sec:at-uncertainty}).

\subsection{Binned counts: a discrete-time approximation}\label{sec:est-inar}

Online volumes, and some offline sources, are available only as counts per hour or per
day. Let $X_{k,t}=N_k((t-1)\Delta,t\Delta]$ be the count of stream $k$ in bin $t$, where
$\Delta$ is the bin width. The integer-valued autoregression INAR($p$) with Poisson
offspring and Poisson immigration is
\begin{equation}\label{eq:inar}
\begin{gathered}
  X_{k,t}\mid\Hist_{(t-1)\Delta}\ \sim\ \Pois\Big(\nu_{k,t}+\sum_{j=1}^M\sum_{\ell=1}^p\varphi_{kj,\ell}\,X_{j,t-\ell}\Big),\\
  \nu_{k,t}=\int_{(t-1)\Delta}^{t\Delta}\Big\{\mu_k(u)+\upsilon_k\sum_eS_e\psi_k(u-\tau_e)\Big\}\,\dd u,
\end{gathered}
\end{equation}
independently over $k$. In the INAR representation, each of the $X_{j,t-\ell}$ incidents
has $\Pois(\varphi_{kj,\ell})$ offspring in bin $t$. Sums of independent Poisson
variables are Poisson, so \eqref{eq:inar} is the identity-link Poisson autoregression
\citep{fokianos2009}. Its conditional likelihood is exact for the binned model. It is concave in the lag
coefficients, and in any linear parametrisation of $\nu$.

Under the parameter mapping $\nu=\Delta\mu$, $\varphi_{kj,\ell}=\Delta\varphi_{kj}(\ell\Delta)$
of \citet{kirchner2016,kirchner2017}, the INAR($\infty$) bin-count process converges to the
Hawkes process as $\Delta\to0$, and conditional least squares (CLS) for the INAR($p$)
representation gives $\hat\alpha_{kj}=\sum_{\ell\le p}\hat\varphi_{kj,\ell}$
(\cref{app:est-inar}). The mapping is a first-order statement. What a fitted
autoregression converges to at a fixed bin width is answered exactly by
\cref{prop:est-binning} in \cref{app:est-inar}. For a stationary univariate Hawkes process
with branching ratio $n$, the coefficients of the best linear predictor of a bin count
from its past, which are the limit of CLS, sum to $n_\Delta=1-(1-n)\,\varpi_\Delta^{1/2}\le n$,
where $\varpi_\Delta\ge1$ is the innovation dispersion; for an exponential kernel,
$n_\Delta=n-\tfrac12n(1-n)\beta\Delta+O(\Delta^2)$. The bias of autoregressive contagion
estimates is therefore of order $\beta\Delta$, and it does not vanish for a fixed bin width,
however long the series: offspring born in their parent's bin are invisible to lagged
regressors. The same identity holds for a Cox process without contagion, so binned
autoregressions are exposed to spurious contagion in the same way as continuous-time fits,
and within-bin cross-excitation, such as same-day online-to-offline spillover, never enters
the lag coefficients (\cref{app:est-inar}).

\paragraph{Nonparametric kernels and sparse excitation.} Histogram and spline kernels
\citep{marsan2008,lewis2011} keep every M-step of \cref{alg:em} in closed form, and
$\ell_1$ or nuclear-norm penalties give sparse or low-rank estimates of $\Amat$ when there
are many streams \citep{hansen2015,zhou2013}. \Cref{sec:est-nonpar} gives the updates and
their costs. Neither is evaluated in this paper.

\subsection{Testing whether one stream excites another}\label{sec:est-lrt}

The hypothesis $H_0:\alpha_{kj}=0$ says that stream $j$ does not excite stream $k$. For
$k\ne j$ it is Granger non-causality from $j$ to $k$ (\cref{prop:granger}). With
$k=\mathrm{off}$ and $j=\mathrm{on}$, it is the hypothesis of no online-to-offline
spillover. Under $H_0$ the parameter lies on the boundary of $\{\alpha_{kj}\ge0\}$, and the
usual $\chi^2_1$ reference is wrong. Write $\vartheta=(\alpha,\vartheta_2)$, where
$\alpha=\alpha_{kj}$ and $\vartheta_2\in\Theta_2\subset\mathbb R^d$ collects the remaining
parameters, and let $\mathrm{LR}_T=2\{\ell_T(\hat\vartheta)-\ell_T(\tilde\vartheta)\}$, where
$\hat\vartheta$ and $\tilde\vartheta$ maximise $\ell_T$ over $[0,\infty)\times\Theta_2$ and
over $\{0\}\times\Theta_2$.

We use the classical result for a parameter on the boundary
\citep{chernoff1954,self1987}, which \cref{app:est-lrt} restates with its conditions and a
proof (\cref{prop:lrt}). The conditions are that both maximisers are consistent, that the
normalised score is asymptotically normal with positive definite information $\bm I$, and
that the observed information converges to $\bm I$ locally uniformly. Under them,
\begin{equation}\label{eq:est-chibar}
  \mathrm{LR}_T\Rightarrow\tfrac12\chi^2_0+\tfrac12\chi^2_1 .
\end{equation}
The asymptotic $p$-value of an observed value $\mathrm{LR}_T=x>0$ is
$\tfrac12\Pr(\chi^2_1\ge x)$, and it is $1$ if $x=0$. The $5\%$ critical value is $2.71$,
not $3.84$. Moreover $\sqrt T\,\hat\alpha$ converges to the positive part of a centred
normal variable, so Wald intervals for $\alpha_{kj}$ are invalid near the boundary.
The conditions require the decay of $\varphi_{kj}$ to be identified under $H_0$: it must be
fixed in advance, or shared with kernels whose branching ratios are positive under $H_0$,
as with a common decay within a block (\cref{sec:plan-models}). \Cref{rem:est-lrt} in
\cref{app:est-lrt} discusses this and the other conditions.

\paragraph{Decays identified only under the alternative.} If $\varphi_{kj}$ has its own
decay rate, or a whole block of entries that share one decay is tested at once, the decay
is not identified under $H_0$ and the chi-bar-square reference is invalid. We then use the
sup-statistic over a grid $\beta_1<\dots<\beta_G$ in a pre-registered interval
$\mathcal B=[\beta_L,\beta_U]$,
\begin{equation}\label{eq:est-suplr}
  \mathrm{LR}_{\sup}=\max_{g\le G}\mathrm{LR}_T(\beta_g),
\end{equation}
where $\mathrm{LR}_T(\beta)$ fixes the decay at $\beta$. Its null limit
\citep{davies1977,davies1987} is stochastically larger than
$\tfrac12\chi^2_0+\tfrac12\chi^2_1$, so our reference is the parametric bootstrap: simulate
the full model at the null fit $\tilde\vartheta$ (\cref{app:computation}), recompute the
sup-statistic with refits, and report
\begin{equation}\label{eq:est-pboot}
  p_{\mathrm B}=\frac{1+\#\{b\le B:\mathrm{LR}^{(b)}_{\sup}\ge\mathrm{LR}_{\sup}\}}{B+1}.
\end{equation}
\Cref{app:est-davies} gives an upcrossing approximation and a cheaper score-based
bootstrap. A common decay within each block (\cref{sec:plan-models}) keeps tests of single
entries within the scope of \eqref{eq:est-chibar} when another entry of the block is
positive, but a test that a whole block vanishes needs the sup-statistic. The analysis plan
registers such tests for its model hierarchy (\cref{sec:plan-models}); none is evaluated in
the simulation study.

\paragraph{Multiple testing.} A full screen of $\Amat$ involves up to $M^2$ tests.
Statistics for different target streams are asymptotically independent when no parameter
is shared across rows, because their scores are stochastic integrals against martingales
without common jumps; statistics for the same target are correlated. The pre-registered
plan controls the family-wise error rate with Holm's procedure \citep{holm1979}
(\cref{sec:plan-models}). For exploratory screens of large matrices we control the false
discovery rate with the step-up procedure of \citet{benjamini1995}, in the variant of
\citet{benjamini2001} that is valid under arbitrary dependence, because scores within a row
can be negatively correlated.

\subsection{Goodness of fit by time rescaling}\label{sec:est-gof}

By the time-rescaling theorem \citep{papangelou1972,brown2002}, restated with a proof as
\cref{prop:rescaling} in \cref{app:est-gof}, if the model is correct then the rescaled
times $\Lam_k(t^k_1)<\Lam_k(t^k_2)<\cdots$ of each stream form a unit-rate Poisson process.
The increments $z^k_i=\Lam_k(t^k_i)-\Lam_k(t^k_{i-1})$ are i.i.d.\ $\mathrm{Exp}(1)$, and the
$u^k_i=1-e^{-z^k_i}$ are i.i.d.\ uniform on $(0,1)$. The positivity condition holds because
$\lambda_k\ge\mu_k>0$. The points of different streams, rescaled by their own compensators,
are moreover independent \citep{meyer1971,brown1988}, so increments can be pooled across
streams.

\paragraph{Diagnostics.} For each stream we use a Kolmogorov--Smirnov test of the $u^k_i$
against the uniform law, displayed as a uniform quantile plot with the asymptotic $95\%$
band $\pm1.36/\sqrt{n_k}$ \citep{brown2002}; a lag plot and the autocorrelations of the
$u^k_i$, which check independence; and the residual process $N_k(t)-\Lam_k(t)$ against
$t$ \citep{ogata1988}, which shows \emph{where} in time a model fails. A misspecified
background shows up as long excursions, a misspecified shock response as excursions that
start at the $\tau_e$.

\paragraph{Estimated parameters.} With $\hat\vartheta$ in place of $\vartheta_0$, the
$u^k_i$ are no longer exactly i.i.d.\ uniform, and the plug-in Kolmogorov--Smirnov
$p$-values are not calibrated. In the i.i.d.\ setting with maximum likelihood estimates the
plug-in test is conservative \citep{durbin1973,lilliefors1967}; \cref{app:est-gof} explains
why the same mechanism plausibly operates here and when it can fail. Either way, we
calibrate by parametric bootstrap, as pre-registered in \cref{sec:plan-validation}. Simulate
from $\hat\vartheta$ with the observed calendar, refit, recompute the Kolmogorov--Smirnov
statistic $\mathrm{KS}^{(b)}$, and report $(1+\#\{\mathrm{KS}^{(b)}\ge\mathrm{KS}\})/(B+1)$.
The time-rescaling check is not evaluated in the simulation study.

\paragraph{What the check can and cannot show.} Time rescaling uses the full conditional
intensity, so in principle it sees beyond second order: a Hawkes model fitted to a
shot-noise Cox process misplaces where the intensity is high. In moderate samples,
however, its power against the matched Cox processes of \cref{thm:equiv} may be low. A
passed check supports the fitted model as a description of the data; it does not
establish contagion.

\subsection{Placebo calibration of shock effects}\label{sec:est-placebo}

The loadings $\upsilon_k$ and the direct components $D_k$ of \cref{thm:attribution} are
estimated in a model whose background and kernels are approximations. Suppose the
background is too rigid around the dates of salient events, or that the calendar overlaps
with other periods of elevated activity. Then an extra shock term would pick up a response
at almost any date, and the asymptotic test of $\upsilon_k=0$, itself a boundary test, would
be miscalibrated. Placebo calendars provide a reference distribution that does not rely on
the model being correct. The calibration is done one class $c$ of real events at a time.
The \emph{base model} $\mathcal M_{-c}$ is the selected model with class $c$ removed. A
\emph{design} $\mathcal R_c$ draws calendars with the same number of events as class $c$:
dates uniformly from eligible days, with the sizes $S_e$ of class $c$ randomly permuted over
them. The eligible days and any spacing rule may depend on the window and on the calendars
of the other classes, but not on the real dates of class $c$, and they must include every
day on which a real event could fall; excluding holidays, for instance, rules out the
randomisation hypothesis below whenever a real event falls on a holiday. The design should
also reproduce rules that constrain the real calendar, such as merging events less than 48
hours apart (\cref{sec:plan-events}). Let $C_0$ be the real class-$c$ calendar and
$C_1,\dots,C_B$ independent draws from $\mathcal R_c$. For $b=0,\dots,B$, fit
$\mathcal M_{-c}$ with $C_b$ added as a shock class and record the same \emph{statistic}
$\Gamma^{(b)}=\Gamma(C_b;\text{data})$: the loading, its likelihood-ratio statistic, or the
direct component $\hat D_k$. \emph{Report} the rank $p$-value
$p=(1+\#\{b\ge1:\Gamma^{(b)}\ge\Gamma^{(0)}\})/(B+1)$, and the rejection frequency of the
asymptotic shock test across the placebo fits, which is the empirical false-positive rate
at the nominal level.

\begin{lemma}[Validity of the placebo rank test]\label{lem:est-placebo}
Let $\Gamma(\cdot\,;\text{data})$ be computed by the same procedure for every calendar. Suppose that, conditionally on the incident data and on the calendars of the
other classes, $C_0,C_1,\dots,C_B$ are exchangeable. This holds, for example, when
$C_1,\dots,C_B$ are drawn independently from $\mathcal R_c$ and, under the
\emph{randomisation hypothesis}, $C_0$ is a further independent draw from $\mathcal R_c$
that is conditionally independent of the incidents. Then $\Pr(p\le a)\le a$ for every
$a\in[0,1]$, whatever the misspecification of $\mathcal M_{-c}$.
\end{lemma}

The proof is in \cref{app:est-proofs}.

Two tempting shortcuts break exchangeability and make the rank test liberal: excluding
placebo dates near the real class-$c$ dates, which makes the design depend on $C_0$; and
fitting the placebo classes with the real class still in the model, which makes the
statistics different functions of the data. Under the alternative, placebo dates can fall
inside real response windows. This reduces power and makes the false-positive estimate
conservative, but does not affect validity under the null.

The randomisation hypothesis is strong. Salient events cluster in periods of tension, which
may differ systematically from typical eligible days. Without it, the rank $p$-value is
best read as a calibration, not a test: does the base model's response at the real dates
exceed what the same model attributes to comparable dates drawn without reference to them?
The placebo distribution of $\hat D_k$ then estimates the false-attribution floor, the
number of incidents the model would credit to an arbitrary date. A real-event direct
component is interpreted only if it ranks high among its placebo counterparts, and if the
empirical false-positive rate clearly exceeds the nominal level, inference on shocks uses
the placebo distribution instead of the asymptotic reference
(\cref{sec:plan-validation,sec:plan-decisions}). Placebo calibration is not evaluated in
the simulation study.
